\section{Only finitely many cyclotomic prime towers are needed for geometric descent}

\paragraph{BR-15-PRIME-SUPPORT.}
Work in the rational numerical abelian--Artin--Tate categories and with
the exact transition formulas of CC-AG08--09. Let $\xi$ be a geometric
Weil germ of weight $w$, let $K$ be its stabilized centre, and let
$d_\xi$ be its geometric Schur index. Choose a representative
\[
                  \xi=[\gamma,h],\qquad \gamma\in K^\times,
\]
at the least positive exponent $h=h_\xi$ of CC-AG13. Thus
$\mathbb Q(\gamma^n)=K$ for every $n\ge1$. Fix the source field
$\mathbb F_{p^a}$, and let $H_a(\xi)$ be the additive monoid of
multiplicities of the geometric simple $S_\xi$ obtained by base change
from that field. For a source Weil label $\alpha$ with
$[\alpha,a]=\xi$, put
\[
 L_\alpha=\mathbb Q(\alpha),\qquad
 r_\alpha=[L_\alpha:K]\,d_\alpha/d_\xi.
                                                               \tag{PS.1}
\]
The positive integer $r_\alpha$ is the transition multiplicity already
proved in CC-AG09.

Then $H_a(\xi)$ is generated by the $r_\alpha$ for which
\[
       \alpha^h=\gamma^a\theta,
       \qquad \theta\in\mu_{h\text{-primary}}.
                                                               \tag{PS.2}
\]
Here the subscript means that the order has no prime divisor outside
those of $h$. In fact every source label can be replaced by a label
satisfying (PS.2), whose centre is a subfield of the original centre
and whose transition multiplicity
\emph{divides} the original multiplicity. No finite bound on the
exponents of the primes dividing $h$ is asserted.

The assertion concerns a geometric target. Removing a root-of-unity
factor need not preserve a prescribed finite-field target label, so
no such finite-target assertion is included.

\begin{proof}
Take any source simple contributing to the orbit of $\xi$, and conjugate
its label so that $[\alpha,a]=[\gamma,h]$. There is an $N\ge1$ with
\[
                      \alpha^{hN}=\gamma^{aN}.
                                                               \tag{PS.3}
\]
Consequently $K\subset L=\mathbb Q(\alpha)$, because the right side
generates $K$, and
\[
                         \theta=\alpha^h/\gamma^a\in\mu(L).
\]
Write uniquely $\theta=\theta_h\theta_0$, where the order of
$\theta_h$ has only prime divisors dividing $h$, and the order $m$ of
$\theta_0$ is coprime to $h$. Both factors are powers of $\theta$ and
therefore belong to $L$. Choose an integer $e$ with $he\equiv1\pmod m$
(with $\theta_0=1$ if $m=1$), and put
\[
                \zeta=\theta_0^e,\qquad \alpha'=\alpha/\zeta.
                                                               \tag{PS.4}
\]
Then $\zeta^h=\theta_0$, so $(\alpha')^h=\gamma^a\theta_h$.
Multiplication by a root of unity preserves the Weil absolute values,
weight, and the condition of being a unit away from $p$. Thus $\alpha'$
is a valid label over the same source field, and its geometric germ
is still $\xi$. Let $L'=\mathbb Q(\alpha')$. We have $L'\subset L$,
and some positive power of $(\alpha')^h$ is a positive power of
$\gamma^a$; stabilization gives $K\subset L'$.

Let $d'$ and $d$ be the source Schur indices attached to $\alpha'$
and $\alpha$, respectively. We prove directly from the explicit local
invariants that
\[
                              d'\mid[L:L']d.
                                                               \tag{PS.5}
\]
This avoids replacing actual multiplicities by a virtual-object gcd.
The classification input is precisely the formula in CC-AG08: at a
finite place above $p$ the invariant is the negative normalized slope
times the local degree over $\mathbb Q_p$; at a real place it is $w/2$
modulo $\mathbb Z$; at the remaining places it is zero. The Schur index
is the least common denominator of these finitely many invariants.
Write $b'_v$ and $b_z$ for the two sets of invariants.

If $z$ lies over a $p$-adic place $v$ of $L'$, the quotient
$\alpha/\alpha'$ is a root of unity and has valuation zero. Thus
\[
 \begin{split}
 b_z
   &=-\frac{\operatorname{ord}_v(\alpha')}{
                    \operatorname{ord}_v(p^a)}[L_z:\mathbb Q_p]\\
   &=[L_z:L'_v]b'_v\pmod{\mathbb Z}.
 \end{split}                                                   \tag{PS.6}
\]
At the other finite places the same congruence holds with both sides
zero. At an archimedean place over a complex place it also holds with
both sides zero. For a real place extending a real place the local
degree is one and the weights agree. For a complex place over a real
place the local degree is two, which sends $w/2$ to the integer $w$.
Thus the congruence holds at every place, including the real-centre case.

By the definition of $d$, each $d b_z$ is zero in $\mathbb Q/\mathbb Z$.
Sum (PS.6), multiplied by $d$, over the places $z$ above a fixed $v$.
The local degrees sum to the global degree $[L:L']$: this follows by
tensoring the $L'$-vector space $L$ with its completion and using the
product decomposition into the completions of $L$. Consequently
\[
                              [L:L']d b'_v=0
                           \quad\text{in }\mathbb Q/\mathbb Z.
                                                               \tag{PS.7}
\]
It holds for every $v$. Taking the least common denominator of the
$b'_v$ gives (PS.5). The completion decomposition is Milne (2020),
\emph{Algebraic number theory}, Propositions~8.1--8.2, p.~136;
its real and complex version is the same scalar-extension dimension
identity. No Galois assumption on either field extension is needed.

Applying (PS.1) to the two labels now gives
\[
                 \frac{r_\alpha}{r_{\alpha'}}
                          =\frac{[L:L']d}{d'}\in\mathbb Z_{>0}.
                                                               \tag{PS.8}
\]
Thus every original generator of $H_a(\xi)$ is a positive integer
multiple of one with the restricted root-of-unity support (PS.2).
Conversely every restricted label is an allowed source label, so its
multiplicity already belongs to $H_a(\xi)$. Taking finite positive sums
proves equality of the two monoids.
\end{proof}

\paragraph{Exact remaining gate and source scope.}
The unrestricted geometric search over all root-of-unity orders reduces
to towers at the finitely many primes dividing the computable integer
$h_\xi$. The new statement does not bound the height in those towers,
compute a whole-monoid conductor, or assert that the monoid is cyclic.
CC-AG14 already disproves the last assertion. CC-AG12's separate
rank-dependent finite procedure remains valid for every specified object.
The division-algebra and transition inputs are Milne (1994), corrected
author manuscript of 2 November 2019, pp.~17--20, particularly
Theorems~2.15(c), 2.16 and 2.18--2.19, in the established abelian sector;
the proof above supplies the root-of-unity reduction and actual positive
multiplicity divisibility in full.
